\begin{table}[htbp]\centering
\def\sym#1{\ifmmode^{#1}\else\(^{#1}\)\fi}
\caption{Political Participation and Social Media Usage}
\begin{tabular}{l*{1}{ccccc}}
\toprule
                    &\multicolumn{1}{c}{(1)}&            &            &            &            \\
                    &           N&        Mean&          SD&         Min&         Max\\
\midrule
\textit{Outcome Variable}&            &            &            &            &            \\
Voted in Presidential Election&        8947&        0.99&        0.11&           0&           1\\
\textit{Social Media Usage}&            &            &            &            &            \\
Uses Facebook       &        8947&        0.75&        0.43&           0&           1\\
Uses X (Twitter)    &        8947&        0.31&        0.46&           0&           1\\
Facebook Usage Frequency (0-3)&        8947&        1.91&        1.26&           0&           3\\
X Usage Frequency (0-3)&        8947&        0.58&        0.99&           0&           3\\
\bottomrule
\multicolumn{6}{l}{\footnotesize Sample: 9,093 respondents from ANES 2020 and 2024.}\\
\multicolumn{6}{l}{\footnotesize Voted and platform usage (fb\\_use, x\\_use) are binary indicators.}\\
\multicolumn{6}{l}{\footnotesize 98.8\% of sample voted (N=8,985), compared to 66\% national turnout.}\\
\multicolumn{6}{l}{\footnotesize 75.1\% use Facebook; 30.8\% use X.}\\
\multicolumn{6}{l}{\footnotesize Usage frequency categories: 0=not at all, 1=infrequent (less than weekly),}\\
\multicolumn{6}{l}{\footnotesize 2=weekly (1-4 times/week), 3=daily or more (once+ per day).}\\
\multicolumn{6}{l}{\footnotesize Mean FB frequency (1.90) = between infrequent and weekly.}\\
\multicolumn{6}{l}{\footnotesize Mean X frequency (0.58) = between non-user and infrequent.}\\
\end{tabular}
\end{table}
